\chapter{部分问题解答} \label{Problems}

\noindent\textbf{问题} \ref{Pro:scalar_with_uncertainty} \textbf{解：}

仍定义参考模型如式 \eqref{Sys:MRAC:scalar:ref}。
提出下述控制律：
\[ u =  \hat{k}_1 (t) x + \hat{k}_2 (t) u_c - \varphi (x) \hat{\theta}_1 (t) \]
代入 \eqref{Sys:MRAC:nonlinear} 得到闭环系统方程
\begin{equation*}
    \dot{x} = a  x +  b \hat{k}_1 (t) x + b \hat{k}_2 (t) u_c -  \varphi  (x) (b\hat{\theta}_1 (t) - \theta) 
\end{equation*}
类似上一小节对于理想情形的推导，可得与 \eqref{k1*k2*} 一样的$k^{\ast}_1, k^{\ast}_2$。仍定义$e=x-x_{\operatorname{ref}}$，结合 \eqref{Sys:MRAC:scalar:ref} 与上式，误差动力学可写为
\begin{align}
  \dot{e} & =  a  x +  b \hat{k}_1 (t) x + b \hat{k}_2 (t) u_c -\textcolor{second}{  \varphi  (x) (b\hat{\theta}_1 (t) - \theta) }+ a_{\refval}
  x_{\refval} - b_{\refval}
  u_c \nonumber\\
  & = (a_{\refval} + a) x +  b \hat{k}_1 (t) x +
  (b \hat{k}_2 (t) - b_{\refval}) u_c -  \textcolor{second}{\varphi (x)
  (b\hat{\theta}_1 (t) - \theta)} - a_{\refval} (x-x_{\operatorname{ref}})
  \nonumber\\
  & =  b (\hat{k}_1 (t) - k^{\ast}_1) x + b (\hat{k}_2 (t) - k^{\ast}_2)
  u_c -  \textcolor{second}{\varphi (x) (b\hat{\theta}_1 (t) - \theta)} -
  a_{\refval} e \nonumber\\
  & =  - a_{\refval} e + b \tilde{k}_1 x + b
  \tilde{k}_2 u_c - \textcolor{second}{b \varphi \tilde{\theta}} \label{Pro:err:MRAC:nonlinear}
\end{align}
其中定义了三个参数估计误差 $\tilde{k}_1 \triangleq \hat{k}_1 (t) - k^{\ast}_1$， $\tilde{k}_2 \triangleq \hat{k}_2
(t) - k^{\ast}_2$ 和 \textcolor{second}{$\tilde{\theta} \triangleq\hat{\theta}_1 (t) - \frac{\theta}{b}$}。

同样考察下述候选Lyapunov函数
\begin{equation*}
  V = \frac{1}{2} e^2 + \frac{| b |}{2 \gamma_1} \tilde{k}^2_1 + \frac{| b
  |}{2 \gamma_2} \tilde{k}^2_2 + \frac{| b |}{2 \gamma_3} \tilde{\theta}^2 
\end{equation*}
其中 $\gamma_1, \gamma_2, \gamma_3 > 0$.

沿 \eqref{Pro:err:MRAC:nonlinear} 的状态轨线，求其对时间的导数得
\begin{equation*}
    \begin{aligned}\dot{V}&= e\dot{e}+\frac{|b|}{\gamma_{1}}\tilde{k}_{1}\dot{\hat{k}}_{1}+\frac{|b|}{\gamma_{2}}\tilde{k}_{2}\dot{\hat{k}}_{2}+\frac{|b|}{\gamma_{3}}\tilde{\theta}\dot{\hat{\theta}}_{1}\\
    &=-a_{\mathrm{ref}}e^{2}+b\tilde{k}_{1}xe+b\tilde{k}_{2}u_{c}e-b\varphi\tilde{\theta}e+\frac{|b|}{\gamma_{1}}\tilde{k}_{1}\dot{\hat{k}}_{1}+\frac{|b|}{\gamma_{2}}\tilde{k}_{2}\dot{\hat{k}}_{2}+\frac{|b|}{\gamma_{3}}\tilde{\theta}\dot{\hat{\theta}}_{1}\\
    &=-a_{\mathrm{ref}}e^{2}+\frac{|b|}{\gamma_{1}}\tilde{k}_{1}\left(\dot{\hat{k}}_{1}+\gamma_{1}\mathrm{sgn}(b)xe\right)+\frac{|b|}{\gamma_{2}}\tilde{k}_{2}\left(\dot{\hat{k}}_{2}+\gamma_{2}\mathrm{sgn}(b)u_{c}e\right)+\frac{|b|}{\gamma_{3}}\tilde{\theta}\left(\dot{\hat{\theta}}_{1}-\gamma_{3}\mathrm{sgn}(b)\varphi e\right)\end{aligned}
\end{equation*}
则选取
\begin{equation*}
  \dot{\hat{k}}_1 = -  \gamma_1 \sgn (b) x  e,
  \dot{\hat{k}}_2 = -  \gamma_2 \sgn (b)  u_c e,
  \dot{\hat{\theta}}_1 = \gamma_3 \sgn (b)   \varphi e
\end{equation*}
即得
\begin{equation*}
  \dot{V} = - a_{\refval} e^2 \leq 0
\end{equation*}
剩余步骤同此问题之前的例子。

\noindent\textbf{问题} \ref{Pro:uncertainty} \textbf{解：}

假设满足 \eqref{matching:condition:GLS} 的$K^{\ast}_1, K^{\ast}_2$存在。设计控制律\[ u = \hat{K}_1 (t) x + \hat{K}_2 (t) u_c - \hat{\Theta}^\mathrm{T}\Phi (x)  \]
代入 \eqref{Sys:MRAC:MIMO_with_dis} 可得
\begin{align}
    \dot{x}&=A  x + B \Lambda (\hat{K}_1 (t) x + \hat{K}_2 (t) u_c-\textcolor{second}{\hat{\Theta}^\mathrm{T}\Phi})+\textcolor{second}{B\Lambda\Theta^\mathrm{T}\Phi}+
    A_{\refval} x +
    B_{\refval} u_c-A_{\refval} x -
    B_{\refval} u_c\nonumber\\
    &=A_{\refval} x +
    B_{\refval} u_c+(A + B \Lambda \hat{K}_1 (t)-A_{\refval}) x+(B\Lambda\hat{K}_2 (t)-B_{\refval}) u_c-\textcolor{second}{(B\Lambda\hat{\Theta}^\mathrm{T}\Phi-B\Lambda\Theta^\mathrm{T}\Phi)}\nonumber\\
    &=A_{\refval}  x +
   B_{\refval} u_c + B \Lambda  \tilde{K}_1 x + B
   \Lambda  \tilde{K}_2 u_c -\textcolor{second}{B\Lambda\tilde{\Theta}^\mathrm{T}\Phi}\label{mimo_with_dis_dotx}
\end{align}
其中的参数估计误差 $\tilde{K}_1 \triangleq \hat{K}_1 {- K_1^{\ast}} , \tilde{K}_2 \triangleq \hat{K}_2 {- K_2^{\ast}} ,\tilde{\Theta}\triangleq\hat{\Theta}-\Theta$。

定义跟踪误差 $e \triangleq x - x_{\refval}$。根据 \eqref{ideal_mimo} 和 \eqref{mimo_with_dis_dotx}，写出误差动力学
\begin{equation}
  \dot{e} = A_{\refval} e + B \Lambda  \tilde{K}_1 x +
  B \Lambda  \tilde{K}_2 u_c-\textcolor{second}{B\Lambda \tilde{\Theta}^\mathrm{T}\Phi} \label{errsys_with_dis:MRAC:GLS}
\end{equation}

考虑如下候选Lyapunov函数
\begin{equation}
  V=e^\mathrm{T}Pe+\tr(\tilde{K}_1^\mathrm{T}{\Lambda}\tilde{K}_1)+\tr(\tilde{K}_2^\mathrm{T}{\Lambda}\tilde{K}_2)+\textcolor{second}{\tr (\tilde{\Theta}{\Lambda}\tilde{\Theta}^\mathrm{T})}
\end{equation}

求其导数得
\begin{align}
  \dot{V} & =e^\mathrm{T}P\dot{e}+\dot{e}^\mathrm{T}Pe + 2\tr \{
  \tilde{K}^\mathrm{T}_1 \Lambda\dot{\hat{K}} _1 \} + 2\tr
  \{ \tilde{K}^\mathrm{T}_2 \Lambda\dot{\hat{K}} _2 \}+\textcolor{second}{2\tr \{ \tilde{\Theta} \Lambda\dot{\hat{\Theta}}^\mathrm{T} \}}\nonumber\\
  & = e^\mathrm{T}P( A_{\refval} e + B \Lambda  \tilde{K}_1 x +
  B \Lambda  \tilde{K}_2 u_c-\textcolor{second}{B\Lambda \tilde{\Theta}^\mathrm{T}\Phi})+(A_{\refval} e + B \Lambda  \tilde{K}_1 x +
  B \Lambda  \tilde{K}_2 u_c-\textcolor{second}{B\Lambda \tilde{\Theta}^\mathrm{T}\Phi})^\mathrm{T}Pe\nonumber\\
  &\quad+
  2\tr \{ \tilde{K}^\mathrm{T}_1 \Lambda\dot{\hat{K}} _1 \}
  + 2\tr \{ \tilde{K}^\mathrm{T}_2 \Lambda\dot{\hat{K}} _2 \}+\textcolor{second}{2\tr \{ \tilde{\Theta} \Lambda\dot{\hat{\Theta}}^\mathrm{T} \}}\nonumber\\
  &= e^\mathrm{T}(A_{\refval}^\mathrm{T} P + P A_{\refval}) e+ 2 e^\mathrm{T}  P  B \Lambda 
  \tilde{K}_1 x + 2 e^\mathrm{T} P  B \Lambda  \tilde{K}_2 u_c -\textcolor{second}{2 e^\mathrm{T}PB\Lambda \tilde{\Theta}^\mathrm{T}\Phi}\nonumber\\
  &\quad +  2\tr \{ \tilde{K}^\mathrm{T}_1  \Lambda\dot{\hat{K}} _1 \}
  + 2\tr \{ \tilde{K}^\mathrm{T}_2  \Lambda\dot{\hat{K}} _2
  \}+\textcolor{second}{2\tr \{ \tilde{\Theta} \Lambda\dot{\hat{\Theta}}^\mathrm{T} \}}\nonumber\\
  & = -e^\mathrm{T}Q e +
  2\tr \{ \tilde{K}^\mathrm{T}_1 \Lambda B^\mathrm{T}  P  e  x^\mathrm{T} \} +
  2\tr \{ \tilde{K}^\mathrm{T}_2 \Lambda B^\mathrm{T}  P  e  u_c^\mathrm{T} \} - \textcolor{second}{2\tr \{ \tilde{\Theta} \Lambda B^\mathrm{T}Pe\Phi^\mathrm{T}
  \}}\nonumber\\
  &\quad + 2\tr \{ \tilde{K}^\mathrm{T}_1  \Lambda\dot{\hat{K}} _1 \}
  + 2\tr \{ \tilde{K}^\mathrm{T}_2  \Lambda \dot{\hat{K}} _2\}+\textcolor{second}{2\tr \{ \tilde{\Theta} \Lambda\dot{\hat{\Theta}}^\mathrm{T} \}}\label{pro_dot_V_modify}
\end{align}
其中最后一行的转化参看 \eqref{trace_transform}。

选取
\[\dot{\hat{K}} _1 = -  B^\mathrm{T}  P  e  x^\mathrm{T}, \dot{\hat{K}} _2 = -  B^\mathrm{T}  P  e  u^\mathrm{T}_c,\dot{\hat{\Theta}}=\Phi e^\mathrm{T}PB\]
即可消去 \eqref{pro_dot_V_modify} 最后六项。这样就有
\[ \dot{V} = - e^\mathrm{T} Q  e \leq 0 \]
后续分析同 \ref{4Cref} 节。

\noindent\textbf{问题} \ref{Pro:dis_with_unknown_upperbound} \textbf{解：}

设计滑模面$s = x_2 + \lambda x_1 = 0,  \lambda > 0$，$s$的动力学是 $\dot{s} = \dot{x}_2 + \lambda \dot{x}_1 = u + d (t)
  + \lambda x_2$。

定义$\Phi(x)=1 + \| x_1 \| + \| x_2 \|^2\ge 1$。可以设计（其中$\hat{k}$是对$k$的估计。注意，只能估计$d_{\max}$（常值），而不能估计含各状态的乘积项；最后一项的作用类似于之前$\tilde{k}$中的$\eta$）
  \[ u = - \lambda x_2 - \hat{k}\Phi(x) \sgn (s) -\eta\sgn(s)\]
这样，$s$的动力学就变为$\dot{s}=d(t)-\hat{k}\Phi(x) \sgn (s)-\eta\sgn(s)$。

定义$\tilde{k}=\hat{k}-d_{\max}$。考虑如下候选Lyapunov函数
  \[ V = \frac{1}{2} s^2 + \frac{1}{2 \gamma} \tilde{k}^2, \gamma > 0, \]
  求其导数得（注意第三行凑出误差$\tilde{k}$的形式）
  \begin{align*}
    \dot{V} & = s  \dot{s} + \frac{1}{\gamma} \tilde{k} \dot{\hat{k}}\\
    &=s(d(t)-\hat{k}\Phi(x)\sgn(s)-\eta\sgn(s))+\frac{1}{\gamma} \tilde{k} \dot{\hat{k}}\\
    & \leq d_{\max}\Phi(x)|s|-\hat{k}\Phi(x)|s|-\eta|s|+\frac{1}{\gamma} \tilde{k} \dot{\hat{k}}\\
    & = -|s|\Phi(x)\tilde{k}+\frac{1}{\gamma} \tilde{k} \dot{\hat{k}}-\eta|s|\\
    &= \frac{1}{\gamma}\tilde{k} (\dot{\hat{k}}-\gamma |s|\Phi(x))-\eta|s|
  \end{align*}
则取$\dot{\hat{k}}=\gamma |s|\Phi(x)$，就有$\dot{V}\leq-\eta|s|$，之后的分析同 \ref{5C_disturbance} 节其他各例。

\noindent\textbf{问题} \ref{Pro:backstep-complex} \textbf{解：}

考虑下述系统 \begin{align*}
    \dot{x}_1 & = -x_1^3+ x_2+\sin x_1\\
    \dot{x}_2 & = \cos x_2-x_2^2 + x_3 \\
    \dot{x}_3 & = u +\cos x_1 +x_2^2
  \end{align*}

\noindent{\textbf{步骤 1: }}

定义 $z_1 = x_1$，将$x_2$看作控制输入。设计“理想的”$x_2$
\begin{align*}
  \Phi_1 (x_1) & = -\sin x_1
\end{align*}
并定义状态变换（“真实的”$x_2$减去“理想的”$x_2$）
\[ z_2 \triangleq x_2 - \Phi_1 (x_1)= x_2 +\sin x_1\]
于是$z_1$的动力学写为
\begin{align*}
  \dot{z}_1 & = \dot{x}_1\\
  & = -x_1^3+ x_2+\sin x_1\\
  & =  -z_1^3+z_2
\end{align*}
考虑如下候选 Lyapunov 函数
\[ V_1 (z_1) = \frac{1}{2} z^2_1 \]
求其导数得到
\begin{align*}
  \dot{V}_1 & = z_1 \dot{z}_1 = -z_1^4+z_1z_2
\end{align*}

\noindent{\textbf{步骤 2: }}

$z_2$的动力学写为
  \begin{align*}
  \dot{z}_2 & = \dot{x}_2 +\dot{x}_1\cos x_1 \\
  & = \cos x_2-x_2^2 + x_3 + (-x_1^3+ x_2+\sin x_1)\cos x_1
\end{align*}

考察下列候选合并Lyapunov函数
\[ V_2 (z_1, z_2) = V_1 (z_1) + \frac{1}{2} z^2_2 \]
求其导数得
\begin{align*}
  \dot{V}_2 & = \dot{V}_1 + z_2 \dot{z}_2\\
  & = -z_1^4+z_1z_2+z_2(\cos x_2-x_2^2 + x_3 + (-x_1^3+ x_2+\sin x_1)\cos x_1)
\end{align*}
则可设计“理想的”$x_3$
\begin{align*}
    \Phi_2(x_1,x_2) &= -z_1-z_2-\cos x_2+x_2^2- (-x_1^3+ x_2+\sin x_1)\cos x_1 \\&= -x_1-x_2-\sin x_1 -\cos x_2+x_2^2-(-x_1^3+ x_2+\sin x_1)\cos x_1 
\end{align*}

定义状态变换
\[z_3 = x_3 - \Phi_2(x_1,x_2) \]
于是可得
\[ \dot{V}_2 = -z_1^4+z_1z_2+z_2(z_3-z_1-z_2) = - z_1^4 - z^2_2+z_2z_3 \]

\noindent{\textbf{步骤 3: }}

$z_3$的动力学写为
  \begin{align*}
  \dot{z}_3 & = \dot{x}_3 - \frac{\partial \Phi_2(x_1,x_2)}{\partial x_1}\dot{x}_1 - \frac{\partial \Phi_2(x_1,x_2)}{\partial x_2}\dot{x}_2
\end{align*}
考虑合并Lyapunov函数
\begin{align*}
  V_c & = V_2  + \frac{1}{2} z^2_3 
\end{align*}
根据上述控制律，其导数为
\begin{align*}
    \dot{V}_c &= - z_1^4 - z^2_2+z_2z_3 +z_3\left(\dot{x}_3 - \frac{\partial \Phi_2(x_1,x_2)}{\partial x_1}\dot{x}_1 - \frac{\partial \Phi_2(x_1,x_2)}{\partial x_2}\dot{x}_2\right)\\
    &=- z_1^4 - z^2_2+z_2z_3 +z_3\left(u+\cos x_1+x_2^2 - \frac{\partial \Phi_2(x_1,x_2)}{\partial x_1}\dot{x}_1 - \frac{\partial \Phi_2(x_1,x_2)}{\partial x_2}\dot{x}_2\right)
\end{align*}
这样，选取\begin{align*}
    u&=-z_2-z_3-\cos x_1-x_2^2+\frac{\partial \Phi_2(x_1,x_2)}{\partial x_1}\dot{x}_1 + \frac{\partial \Phi_2(x_1,x_2)}{\partial x_2}\dot{x}_2\\
    &=-z_2-z_3-\cos x_1-x_2^2+\frac{\partial \Phi_2(x_1,x_2)}{\partial x_1}(-x_1^3+ x_2+\sin x_1) + \frac{\partial \Phi_2(x_1,x_2)}{\partial x_2}(\cos x_2-x_2^2 + x_3)
\end{align*}
即可得到$\dot{V}_c = - z_1^4 - z^2_2-z_3^2$，其为负定。接下来类似该问题之前的推导即可。

\noindent\textbf{问题} \ref{Pro:pos_ctrl_with_Dq} \textbf{解：}

定义位置误差和速度误差：
\begin{align*}
    \tilde{q}(t) \triangleq q(t) - q_d, \quad \dot{\tilde{q}}(t) \triangleq \dot{q}(t) - \dot{q}_d = \dot{q}(t) \text{。}
\end{align*}
从而 $\ddot{\tilde{q}}(t) = \ddot{q}(t)$。为简化记号，在不引起歧义的前提下，可能省略后续一些函数中的参数。

误差动力学为：
\begin{equation*}
    M (q) \ddot{\tilde{q}} + C (q, \dot{q}) \dot{\tilde{q}} + \textcolor{second}{D\dot{\tilde{q}}} + g(q) = \tau \text{，}
\end{equation*}

可采用如下控制输入：
\begin{equation}
    \tau = g(q) - K_p \tilde{q} - K_d \dot{\tilde{q}}, \quad K_p, K_d > 0 \text{，}
\end{equation}
从而闭环系统为：
\begin{equation}
   M (q) \ddot{\tilde{q}} + C (q, \dot{q}) \dot{\tilde{q}} + \textcolor{second}{D\dot{\tilde{q}}} = - K_p \tilde{q} - K_d \dot{\tilde{q}} \text{。}
\end{equation}
考虑候选 Lyapunov 函数
\[
    V = \frac{1}{2} \dot{\tilde{q}}^\mathrm{T} M \dot{\tilde{q}}+\frac{1}{2} \tilde{q}^\mathrm{T} K_p \tilde{q}  \text{，}
\]
其沿系统轨线的导数
\begin{align*}
    \dot{V} & = \dot{\tilde{q}}^\mathrm{T} M \ddot{\tilde{q}} + \frac{1}{2}
    \dot{\tilde{q}}^\mathrm{T} \dot{M} \dot{\tilde{q}} + \dot{\tilde{q}}^\mathrm{T} K_p \tilde{q}\\
    & = \dot{\tilde{q}}^\mathrm{T} (- K_p \tilde{q} - K_d \dot{\tilde{q}} - C
    \dot{\tilde{q}}-\textcolor{second}{D\dot{\tilde{q}}}) + \frac{1}{2} \dot{\tilde{q}}^\mathrm{T} \dot{M} \dot{\tilde{q}}+ \dot{\tilde{q}}^\mathrm{T} K_p \tilde{q}\\
    & =  - \dot{\tilde{q}}^\mathrm{T} K_d
    \dot{\tilde{q}} + \frac{1}{2} \dot{\tilde{q}}^\mathrm{T} (\dot{M} - 2 C)
    \dot{\tilde{q}}- \textcolor{second}{\dot{\tilde{q}}^\mathrm{T}D\dot{\tilde{q}}}\\
    & =  - \dot{\tilde{q}}^\mathrm{T} K_d
    \dot{\tilde{q}} - \textcolor{second}{\dot{\tilde{q}}^\mathrm{T}D\dot{\tilde{q}}}
\end{align*}
从而，$V$ 是正定的，$\dot{V}$是负半定的。接下来类似 \ref{6Bref} 节利用 LaSalle 不变集原理（\ref{lasalle}），即可说明原点是渐近稳定的，即 $\lim\limits_{t \to \infty} q (t) = q_d$ 且 $\lim\limits_{t \to \infty} \dot{q}(t) = 0$。

\noindent\textbf{问题} \ref{Pro:robot_disturbance_known} \textbf{解：}

沿用 \ref{6Cref}、\ref{6Dref} 节的记号，$q_d,\dot{q}_d,\ddot{q}_d$ 均有界，且
\[
    s=\dot{\tilde{q}}+\lambda\tilde{q},\quad
    \dot{q}_r=\dot{q}_d-\lambda\tilde{q},\quad \lambda>0 \text{。}
\]
下面分别采用普通滑模方法与平滑方法，均沿用 \ref{5C_disturbance} 节中的设计。令
\[
    \tau_0=g(q)+M\ddot{q}_r+C\dot{q}_r,\quad
    \tau=\tau_0-Ks-v,\quad K=K^\mathrm{T}>0 \text{，}
\]
其中 $v$ 是待设计的扰动补偿项。闭环系统为
\[
    M\dot{s}+Cs=-Ks-v+d(t) \text{。}
\]
记 $V_s=\frac12s^\mathrm{T}Ms$。由 $\dot{M}-2C$ 的反对称性，有
\[
    \dot{V}_s=-s^\mathrm{T}Ks-s^\mathrm{T}v+s^\mathrm{T}d(t) \text{。}
\]
以下符号函数按分量作用，且 $s^\mathrm{T}\sgn(s)=\sum_{i=1}^p|s_i|\triangleq\|s\|_1\geq\|s\|$；在零点处采用理想滑模约定，开关值允许在 $[-1,1]$ 内取值。平滑方法中均取 $\varepsilon(t)=\mathrm{e}^{-t}$。

\noindent\textbf{（1）普通滑模方法。}取
\[
    v=k\sgn(s),\qquad k>d_{\mathrm{max}} \text{。}
\]
由 $s^\mathrm{T}d(t)\leq d_{\mathrm{max}}\|s\|\leq d_{\mathrm{max}}\|s\|_1$，可得
\begin{align*}
    \dot{V}_s
    &\leq -s^\mathrm{T}Ks-(k-d_{\mathrm{max}})\|s\|_1\\
    &\leq -\lambda_{\min}(K)\|s\|^2
    \leq -\frac{2\lambda_{\min}(K)}{k_{\overline{m}}}V_s \text{。}
\end{align*}
故 $V_s$ 指数衰减，$s(t)\to0$。

\noindent\textbf{（2）平滑方法。}取
\[
    v=\frac{k^2}{k\|s\|+\varepsilon(t)}s,\qquad k>0,\quad k\geq d_{\mathrm{max}} \text{。}
\]
此时
\begin{align*}
    \dot{V}_s
    &= -s^\mathrm{T}Ks-\frac{k^2\|s\|^2}{k\|s\|+\varepsilon(t)}+s^\mathrm{T}d(t)\\
    &\leq -s^\mathrm{T}Ks-\frac{k^2\|s\|^2}{k\|s\|+\varepsilon(t)}+k\|s\|\\
    &= -s^\mathrm{T}Ks+\frac{\varepsilon(t)k\|s\|}{k\|s\|+\varepsilon(t)}\\
    &\leq -\lambda_{\min}(K)\|s\|^2+\mathrm{e}^{-t} \text{。}
\end{align*}
两边积分得到
\[
    V_s(t)+\lambda_{\min}(K)\int_0^t\|s(\xi)\|^2\diff\xi
    \leq V_s(0)+1 \text{，}
\]
故 $s\in\mathbb{L}_2\cap\mathbb{L}_\infty$。

两种方法中，均有 $s\in\mathbb{L}_2\cap\mathbb{L}_\infty$。由 $\dot{\tilde{q}}=-\lambda\tilde{q}+s$ 的输入-状态稳定性，$\tilde{q},\dot{\tilde{q}},q,\dot{q}$ 均有界。普通方法中的符号函数有界，平滑方法中的补偿项满足
\[
    \left\|\frac{k^2s}{k\|s\|+\varepsilon(t)}\right\|\leq k \text{，}
\]
因此 $v$ 有界。结合 $M^{-1}$ 的有界性及 $\|Cs\|\leq k_c\|\dot{q}\|\|s\|$，由闭环方程可得 $\dot{s}$ 有界。控制输入也有界，解对所有 $t\geq0$ 存在。由于 $\|s\|^2$ 一致连续且其积分有界，由 Barbalat 引理可得 $s(t)\to0$。再由输入-状态稳定性，$\tilde{q}(t)\to0$，并由 $\dot{\tilde{q}}=-\lambda\tilde{q}+s$ 得 $\dot{\tilde{q}}(t)\to0$，故两种设计均实现跟踪控制。

\noindent\textbf{问题} \ref{Pro:robot_disturbance_unknown} \textbf{解：}

沿用问题 \ref{Pro:robot_disturbance_known} 解中的记号，仍有 $\|d(t)\|\leq d_{\mathrm{max}}$，但常数 $d_{\mathrm{max}}$ 未知。
将未知常数 $d_{\mathrm{max}}$ 的估计记为 $\hat{k}$，定义 $\tilde{k}=\hat{k}-d_{\mathrm{max}}$。取 $\gamma>0$、$\hat{k}(0)\geq0$，考虑
\[
    V=V_s+\frac{1}{2\gamma}\tilde{k}^2 \text{。}
\]

\noindent\textbf{（1）普通自适应滑模方法。}设计
\[
    v=\hat{k}\sgn(s),\qquad \dot{\hat{k}}=\gamma\|s\|_1 \text{。}
\]
于是
\begin{align*}
    \dot{V}
    &=-s^\mathrm{T}Ks-\hat{k}\|s\|_1+s^\mathrm{T}d(t)
    +\frac{1}{\gamma}\tilde{k}\dot{\hat{k}}\\
    &\leq -s^\mathrm{T}Ks-\tilde{k}\|s\|_1
    +\frac{1}{\gamma}\tilde{k}\dot{\hat{k}}\\
    &=-s^\mathrm{T}Ks \text{。}
\end{align*}
两边积分可得
\[
    V(t)+\lambda_{\min}(K)\int_0^t\|s(\xi)\|^2\diff\xi\leq V(0) \text{，}
\]
故 $s\in\mathbb{L}_2\cap\mathbb{L}_\infty$，且 $\tilde{k},\hat{k}$ 有界。

\noindent\textbf{（2）平滑自适应方法。}设计
\[
    v=\frac{\hat{k}^2}{\hat{k}\|s\|+\varepsilon(t)}s,\qquad
    \dot{\hat{k}}=\gamma\|s\| \text{。}
\]
因 $\dot{\hat{k}}\geq0$，始终有 $\hat{k}\geq0$。由扰动界及自适应律，得到
\begin{align*}
    \dot{V}
    &\leq -s^\mathrm{T}Ks-\frac{\hat{k}^2\|s\|^2}{\hat{k}\|s\|+\varepsilon(t)}
    +d_{\mathrm{max}}\|s\|+\tilde{k}\|s\|\\
    &=-s^\mathrm{T}Ks+\frac{\varepsilon(t)\hat{k}\|s\|}{\hat{k}\|s\|+\varepsilon(t)}\\
    &\leq -\lambda_{\min}(K)\|s\|^2+\mathrm{e}^{-t} \text{。}
\end{align*}
两边积分可得
\[
    V(t)+\lambda_{\min}(K)\int_0^t\|s(\xi)\|^2\diff\xi\leq V(0)+1 \text{，}
\]
故同样有 $s\in\mathbb{L}_2\cap\mathbb{L}_\infty$，且 $\tilde{k},\hat{k}$ 有界。


上述两种设计中，$s$ 均有界。由 $\dot{\tilde{q}}=-\lambda\tilde{q}+s$ 的输入-状态稳定性，$\tilde{q},\dot{\tilde{q}},q,\dot{q}$ 均有界。普通方法中的符号函数有界，平滑方法中的补偿项则满足
\[
    \|\hat{k}\sgn(s)\|\leq\sqrt{p}\hat{k},\qquad
    \left\|\frac{\hat{k}^2s}{\hat{k}\|s\|+\varepsilon(t)}\right\|\leq\hat{k} \text{，}
\]
因此 $v$ 有界。结合 $M^{-1}$ 的有界性及 $\|Cs\|\leq k_c\|\dot{q}\|\|s\|$，由闭环方程可得 $\dot{s}$ 有界。控制输入及自适应律中的 $\dot{\hat{k}}$ 也有界，解对所有 $t\geq0$ 存在。由于 $\|s\|^2$ 一致连续且其积分有界，由 Barbalat 引理可得 $s(t)\to0$。再由输入-状态稳定性，$\tilde{q}(t)\to0$，并由 $\dot{\tilde{q}}=-\lambda\tilde{q}+s$ 得 $\dot{\tilde{q}}(t)\to0$，故两种设计均实现跟踪控制。

\noindent\textbf{问题} \ref{Pro:robot_disturbance_state} \textbf{解：}

常数扰动界未知时的两种方法已在上一题给出。对于状态相关的扰动界，记
\[
    \Phi(q,\dot{q})=1+\|q\|^2+\|\dot{q}\|,\qquad
    \|d(t)\|\leq d_{\mathrm{max}}\Phi(q,\dot{q}) \text{。}
\]
其中 $d_{\mathrm{max}}$ 是未知的非负常数，$\Phi$ 是已知函数。沿用上一题的记号和参考轨迹条件，取 $K=K^\mathrm{T}>0$、$\gamma>0$、$\hat{k}(0)\geq0$。定义 $\tilde{k}=\hat{k}-d_{\mathrm{max}}$，两种方法均考虑候选 Lyapunov 函数
\[
    V=\frac12s^\mathrm{T}Ms+\frac{1}{2\gamma}\tilde{k}^2 \text{。}
\]

\noindent\textbf{（1）普通自适应滑模方法。}设计
\begin{align*}
    \tau &= g(q)+M\ddot{q}_r+C\dot{q}_r-Ks-\hat{k}\Phi\sgn(s),\\
    \dot{\hat{k}} &= \gamma\Phi\|s\|_1 \text{。}
\end{align*}
符号函数的约定与上一题相同。闭环系统为
\[
    M\dot{s}+Cs=-Ks-\hat{k}\Phi\sgn(s)+d(t) \text{。}
\]
利用 $\dot{M}-2C$ 的反对称性及扰动界，有
\begin{align*}
    \dot{V}
    &=-s^\mathrm{T}Ks-\hat{k}\Phi\|s\|_1+s^\mathrm{T}d(t)
    +\frac{1}{\gamma}\tilde{k}\dot{\hat{k}}\\
    &\leq -s^\mathrm{T}Ks-\tilde{k}\Phi\|s\|_1
    +\frac{1}{\gamma}\tilde{k}\dot{\hat{k}}\\
    &=-s^\mathrm{T}Ks \text{，}
\end{align*}
其中用到了 $s^\mathrm{T}d(t)\leq d_{\mathrm{max}}\Phi\|s\|\leq d_{\mathrm{max}}\Phi\|s\|_1$。两边积分可得
\[
    V(t)+\lambda_{\min}(K)\int_0^t\|s(\xi)\|^2\diff\xi\leq V(0) \text{。}
\]

\noindent\textbf{（2）平滑自适应方法。}取 $\varepsilon(t)=\mathrm{e}^{-t}$，设计
\begin{align*}
    \tau &= g(q)+M\ddot{q}_r+C\dot{q}_r-Ks
    -\frac{\hat{k}^2\Phi^2}{\hat{k}\Phi\|s\|+\varepsilon(t)}s,\\
    \dot{\hat{k}} &= \gamma\Phi\|s\| \text{。}
\end{align*}
因 $\dot{\hat{k}}\geq0$，始终有 $\hat{k}\geq0$。利用 $\dot{M}-2C$ 的反对称性及扰动界，得到
\begin{align*}
    \dot{V}
    &= -s^\mathrm{T}Ks-\frac{\hat{k}^2\Phi^2\|s\|^2}{\hat{k}\Phi\|s\|+\varepsilon(t)}
    +s^\mathrm{T}d(t)+\frac{1}{\gamma}\tilde{k}\dot{\hat{k}}\\
    &\leq -s^\mathrm{T}Ks-\frac{\hat{k}^2\Phi^2\|s\|^2}{\hat{k}\Phi\|s\|+\varepsilon(t)}
    +d_{\mathrm{max}}\Phi\|s\|+\tilde{k}\Phi\|s\|\\
    &= -s^\mathrm{T}Ks+\frac{\varepsilon(t)\hat{k}\Phi\|s\|}{\hat{k}\Phi\|s\|+\varepsilon(t)}\\
    &\leq -\lambda_{\min}(K)\|s\|^2+\mathrm{e}^{-t} \text{。}
\end{align*}
与上一题一样，两边积分可得
\[
    V(t)+\lambda_{\min}(K)\int_0^t\|s(\xi)\|^2\diff\xi\leq V(0)+1 \text{，}
\]
两种方法中，均有 $s\in\mathbb{L}_2\cap\mathbb{L}_\infty$，且 $\tilde{k},\hat{k}$ 有界。由 $\dot{\tilde{q}}=-\lambda\tilde{q}+s$ 可得 $\tilde{q},\dot{\tilde{q}},q,\dot{q}$ 有界，因而 $\Phi$ 和 $d(t)$ 均有界。两种补偿项分别满足
\[
    \|\hat{k}\Phi\sgn(s)\|\leq\sqrt{p}\hat{k}\Phi,\qquad
    \left\|\frac{\hat{k}^2\Phi^2s}{\hat{k}\Phi\|s\|+\varepsilon(t)}\right\|
    \leq\hat{k}\Phi \text{，}
\]
故补偿项均有界。将相应的补偿项记为 $v$，由闭环方程
\[
    M\dot{s}+Cs=-Ks-v+d(t)
\]
可得 $\dot{s}$ 有界。控制输入及 $\dot{\hat{k}}$ 也有界，解对所有 $t\geq0$ 存在。由于 $\|s\|^2$ 一致连续且其积分有界，由 Barbalat 引理及输入-状态稳定性，两种方法均可得到
\[
    \lim_{t\to\infty}s(t)=0,\qquad
    \lim_{t\to\infty}\tilde{q}(t)=\lim_{t\to\infty}\dot{\tilde{q}}(t)=0 \text{。}
\]
